\begin{scanpage}{0071}{54}
Expressions (9.03) and (9.05) agree with the result obtained by Braunbek [29] for the scalar fringing field in the far zone. It is also interesting to compare these expressions with the precise numerical results obtained by Belkina [34] by the separation of variables method in the spheroidal coordinate system. It turns out that even with $ka=5$ a satisfactory agreement is observed between our approximation method and the rigorous theory. In Figures 20 and 21, graphs of the functions $V^{(1)}(\theta)$ and $V^{(2)}(\theta)$, are presented which allow one to calculate the fringing field on the basis of the equations
\[\left.\begin{aligned}E_\phi=-H_\theta&=\frac{ika^2}2E_{0y}\cdot V^{(1)}(\theta)\frac{e^{ikR}}R\cos\phi,\\E_\theta=H_\phi&=\frac{ika^2}2E_{0y}\cdot V^{(2)}(\theta)\frac{e^{ikR}}R\sin\phi.\end{aligned}\right\}\tag{9.07}\]
The continuous curve corresponds to the rigorous theory [34]. The dash-dot curve corresponds to the field from the uniform part of the current, and the dashed curve corresponds to the field calculated according to Equation (9.03) and (9.05).
\begin{center}\underline{Oblique Irradiation}\end{center}
\sectiontitle{10}{The Physical Optics Approach}
Let us investigate the general case when the plane wave
\[\mathbf E=\mathbf E_0e^{ik(y\sin\gamma+z\cos\gamma)}\tag{10.01}\]
falls on the disk at an arbitrary angle to its axis. Let us take the spherical coordinate system in such a way that the normal to the incident wave front, $\mathbf n$, would lie in the half-plane $\phi=\frac\pi2$ and form an angle $\gamma$ $\left(0<\gamma<\frac\pi2\right)$ with the $z$ axis (Figure 22). Adhering to the investigation procedure used in the previous sections, let us first calculate the fringing field in the physical optics approach.

The uniform part of the current excited on the disk by wave (10.01) is determined by Equation (3.01) and has the components
\[j_x^0=\frac c{2\pi}H_{0y}e^{iky\sin\gamma},\quad j_y^0=-\frac c{2\pi}H_{0x}e^{iky\sin\gamma},\quad j_z^0=0.\tag{10.02}\]
\end{scanpage}
